% Chapter 2

\chapter{Front Matter}% Main chapter title
\label{Chapter2}

\section{Title Page}%
\label{opt:bib}

The first page of the work is the title page. The design of this page is
specified by \citet{web:promotion}. There exist two versions of the title page
for different applications. The first one is the title page for your
submission of the thesis. After a successful defense of your thesis, it needs
to be published. This is done by uploading your work to
\url{http://ul.qucosa.de}. However, this version needs a different title page.
This is the second version of the title page. To create this title page
instead of the submission title page use the class option
\clsopt{bibversion}.\indexdef{bibversion}

The information on this page is provided by some commands in the data file
\file{thesis\_info.tex}. An overview of this commands is shown in
\autoref{tab:titlepage}. Note that most of the information exclusive for the
bib version can only be given after you defense. This is no problem because
you only need this version afterwards.

A little note to the \verb!\setsubtitle!: Not every work has or needs a
subtitle. In this case, you can simply not use the command
anywhere and the class eliminates the appearance of the subtitle everywhere.


\begin{table}
    \captiont{Titlepage commands}{%
        The commands that are used in the preamble of \file{main.tex} to
        input the information of the title page. Also is given if the
        information is used in the submission (sub) version or the bib
        version.
    }%
    \label{tab:titlepage}
    \begin{tabular}{l l l}
        \toprule
        \tblhead{Command} & \tblhead{Description}& \tblhead{Used for}\\
        \midrule
        \verb!\settitle!&The title of the thesis.&sub, bib\\
        \verb!\setsubtitle!&The subtitle of the thesis.&sub, bib\\
        \verb!\setauthortitle!&The academic degree if the author.&sub, bib\\
        \verb!\setauthor!&The full name of the author.&sub, bib\\
        \verb!\setdayofbirth!&The day of birth of the author.&sub, bib\\
        \verb!\setbirthplace!&The birth place of the author.&sub, bib\\
        \verb!\setfilingdate!&The date where you submit your thesis.&sub\\
        \verb!\setfirstreviewer!&The first reviewer of you thesis&bib\\
        \verb!\setsecondreviewer!&The second reviewer of you thesis&bib\\
        \verb!\setdayofdefens!&The day where you defended your thesis&bib\\
        \verb!\setgrade!&The grade that you get on your thesis&bib\\
        \bottomrule
    \end{tabular}
\end{table}


\section{Dedication}%
\label{cmd:dedicate}

The next page after the title page is the Dedication page. Here you can
dedicate your work to someone or make a suitable insider joke. Whatever it is,
it should be short and in the best case one sentence. You can input this
sentence with the \verb!\setdedication! command in the preamble of
\file{main.tex}.

If you do not want to dedicate your work to someone then leave the command
empty (\verb!\setdedication{}!). This eliminates the dedication page. Note
that the content of \verb!\setdedication! can be any latex command so it not
only possible to write text it could also include a picture or a combination
of both.

\section{Bibliographic Data}

The third nonempty page contains the bibliographic data of the thesis. This
is required by \citet{web:promotion}. Most of the data is calculated
automatically. The only exceptions are the keywords of the thesis which
are inserted by the command \verb!\setkeywords! in the preamble of
\file{main.tex}. The keyword should be given as a comma-separated list.

\subsection{Statistics}%
\label{sub:stats}

In the middle of the page are some automatically calculated value. These
values are the core of the bibliographic data. However, they are not mandatory
from the Promotionsordnung \citep{web:promotion}. Thus, you can deactivate
them with the class option \clsopt{nostatistics}\indexdef{nostatistics}.
Then nothing is calculated or shown, and only the keywords and the paper the
thesis is based on are displayed.

\subsection{Thesis Publication}%
\label{sub:baseon}

Most likely you have published already some of the results you present in this
thesis. Such a publication is then a publication on which the thesis is based
on. They can be divided from publications you only cited by the fact that in
some way you retell them in the thesis. This can even mean that long
paragraphs are equal between them and the text of the thesis.  This
information is critical to prevent the accusation of plagiarism. So this
publications belongs to the bibliographic data of the thesis.

To add publications there is the command \verb!\setbase! in the preamble used.
The command gets a comma-separated list of BibTeX keys. Then it automatically
fetches the full information out of your BibTeX file. However, this, of
course, means that this publication must be in the BibTeX file of the thesis.

The look of this paper can be influenced with the class option
\clsopt{baseonauthorstyle}\indexdef{baseonauthorstyle}. This change how the
authors are displayed. More details can be found at \autoref{sub:author}. The
default value is \quo{gf} which means that first the given name comes and then
the family name.

A second possibility to change the look of the author respectively give extra
information of the authors. This is done by annotations. A deeper look into
annotations can be found in \autoref{sub:annotation}. Two things are vary
handy here. First to mark authors that share a first authorship. This make it
possible to mark that you are a first author even when you are not on first
position. The second thing is the possibility to underline specific authors.
This is useful when you want to mark you in the papers. Maybe when you not on
first position or when you have changes your name by a marriage. Even when
both is not given a mark is a good idea to make the reader easy to spot you.

\section{Abstract}%
\label{opt:abstract}

Then the abstract of the work follows. This is limited to four pages
\citep{web:promotion}. The abstract is written in an extra file. By default
this file is placed at \file{frontmatter/abstract.tex}. You can change the
path to the abstract with the class option
\clsopt{abstractpath}\indexdef{abstractpath}. Note that you have to
\emph{omit} the \file{.tex} extension when defining an alternative value for
this option.

A stand-alone version of the abstract is generated automatically. You are
required to hand it in together with your thesis. It will be handed out during
your defense.

\section{German Abstract}%
\label{sec:germanabstract}

The template gives the option to include a German version of your abstract.
This is mandatory when you write a binational thesis. If you only write it for
the University of Leipzig you can choose if you want to include it. However,
most people add a German abstract.

The default path is \file{frontmatter/zusammenfassung} (omit the \file{.tex}
extension). Again you can change the path to the file that contains the German
abstract with the class option \clsopt{zusammenpath}\indexdef{zusammenpath}.
If you do not want the German abstract, you can simply set the class option to
an empty value (\verb!zusammenpath=,!).

\section{Acknowledgments}%
\label{opt:ack}

The third part of the front matter that inputs text from another path is the
acknowledgments. Here you should thank everybody that helped you with the
thesis: family, friends, supervisor, fellow students and maybe me.

The class option \clsopt{acknowledgmentspath}\indexdef{acknowledgmentspath}
changes the path to the file. The default path is
\file{frontmatter/acknowledgments} (omit the \file{.tex} extension).

\section{The Table of Contents}
The \abb{ToC} is directly in front of the main matter. It only contains
entries from the main matter and the back matter. The \abb{ToC} shows entries
down to the section level. Subsections are not shown here. However, they are
shown in mini \abbs{ToC} at the beginning of each chapter. More details on
mini \abbs{ToC} are given in \autoref{sec:chapter}.

% End of chapters/chapter_2.tex
